\chapter{Implementation}

This chapter presents the technical implementation details of the Verdant Search system, demonstrating how the design specifications from Chapter 5 were translated into working code. The implementation spans multiple technology domains including database design with PostgreSQL and pgvector, Python microservices for search and embedding generation, distributed web crawling, and the hybrid BM25+HNSW retrieval pipeline.

\section{Database Implementation}

\subsection{Schema Implementation with PostgreSQL and pgvector}

The database schema was implemented using PostgreSQL 16 with the pgvector extension for efficient vector similarity search. The implementation leverages PostgreSQL's advanced features including JSONB for flexible metadata storage, native array types for position tracking, and GiST/HNSW indices for vector operations.

\lstset{
  language=SQL,
  basicstyle=\ttfamily\small,
  keywordstyle=\color{blue}\bfseries,
  commentstyle=\color{green!60!black}\itshape,
  stringstyle=\color{red},
  showstringspaces=false,
  breaklines=true,
  frame=single,
  numbers=left,
  numberstyle=\tiny\color{gray},
  caption={Database initialization script with BM25 inverted index tables},
  label=code:init_sql
}
\begin{lstlisting}
-- Initialize database with pgvector extension
CREATE EXTENSION IF NOT EXISTS vector;

-- Documents table for storing indexed content
CREATE TABLE IF NOT EXISTS documents (
    id SERIAL PRIMARY KEY,
   title TEXT NOT NULL,
    content TEXT NOT NULL,
    url TEXT,
    source_type VARCHAR(50),
    doc_metadata JSONB,
    doc_length INTEGER DEFAULT 0,  -- Document length (tokens)
    created_at TIMESTAMP DEFAULT CURRENT_TIMESTAMP,
    updated_at TIMESTAMP DEFAULT CURRENT_TIMESTAMP
);

-- Terms table for BM25 inverted index
CREATE TABLE IF NOT EXISTS terms (
    id SERIAL PRIMARY KEY,
    term VARCHAR(255) UNIQUE NOT NULL,
    doc_frequency INTEGER DEFAULT 0,     -- DF: Number of documents containing term
    total_frequency BIGINT DEFAULT 0,    -- Total occurrences across all documents
    created_at TIMESTAMP DEFAULT CURRENT_TIMESTAMP
);

-- Posting List table (Inverted Index)
CREATE TABLE IF NOT EXISTS postings (
    id SERIAL PRIMARY KEY,
    term_id INTEGER NOT NULL REFERENCES terms(id) ON DELETE CASCADE,
    document_id INTEGER NOT NULL REFERENCES documents(id) ON DELETE CASCADE,
    term_frequency INTEGER NOT NULL,     -- TF: Term frequency in document
    positions INTEGER[],                 -- Term positions for phrase queries
    created_at TIMESTAMP DEFAULT CURRENT_TIMESTAMP,
    UNIQUE(term_id, document_id)
);

-- Document statistics for BM25 calculation
CREATE TABLE IF NOT EXISTS doc_stats (
    id SERIAL PRIMARY KEY,
    total_docs INTEGER DEFAULT 0,
    avg_doc_length FLOAT DEFAULT 0,
    updated_at TIMESTAMP DEFAULT CURRENT_TIMESTAMP
);

-- Vector embeddings with HNSW index
CREATE TABLE IF NOT EXISTS document_embeddings (
    id SERIAL PRIMARY KEY,
    document_id INTEGER REFERENCES documents(id) ON DELETE CASCADE,
    embedding vector(512),  -- CLIP embedding dimension
    created_at TIMESTAMP DEFAULT CURRENT_TIMESTAMP
);

-- Create HNSW index for fast ANN search
CREATE INDEX IF NOT EXISTS document_embeddings_hnsw_idx 
ON document_embeddings 
USING hnsw (embedding vector_cosine_ops)
WITH (m = 16, ef_construction = 64);

-- Optimize inverted index with composite indices
CREATE INDEX IF NOT EXISTS postings_term_doc_idx 
ON postings(term_id, document_id);
\end{lstlisting}

The schema design separates the BM25 inverted index (terms, postings, doc\_stats tables) from the vector search infrastructure (document\_embeddings with HNSW index). This modular approach enables independent optimization of each retrieval method while maintaining data integrity through foreign key constraints.

\clearpage

\section{BM25 Search Implementation}

\subsection{SQL-Based BM25 Calculator}

The BM25 retrieval component is implemented using SQL queries that directly operate on the PostgreSQL-resident inverted index. This approach eliminates the need to load the entire index into memory, enabling scalability to millions of documents.

\lstset{
  language=Python,
  basicstyle=\ttfamily\footnotesize,
  keywordstyle=\color{blue}\bfseries,
  commentstyle=\color{green!60!black}\itshape,
  stringstyle=\color{red},
  showstringspaces=false,
  breaklines=true,
  frame=single,
  numbers=left,
  numberstyle=\tiny\color{gray},
  caption={BM25 Calculator implementation with async SQL queries},
  label=code:bm25_calc
}
\begin{lstlisting}
class BM25Calculator:
    """SQL-based BM25 Calculator using database-resident inverted index"""
    
    def __init__(self, k1=1.5, b=0.75):
        self.k1 = k1  # Term saturation parameter
        self.b = b    # Length normalization parameter
        self.tokenizer = get_tokenizer_service()

    async def search(self, query: str, session: AsyncSession, 
                    top_k: int = 10) -> List[Dict[str, Any]]:
        """Perform BM25 search using database-resident inverted index"""
        
        # 1. Tokenize query
        tokens = self.tokenizer.tokenize(query, mode="search")
        if not tokens:
            return []
        
        token_list = list(set(tokens))  # Deduplicate tokens
        
        # 2. Get global statistics
        stats_res = await session.execute(
            text("SELECT total_docs, avg_doc_length FROM doc_stats WHERE id=1")
        )
        stats = stats_res.first()
        
        if not stats or stats.total_docs == 0:
            return []
        
        total_docs, avg_doc_length = stats
        
        # 3. Calculate BM25 scores using SQL
        # Formula: IDF * ((TF * (k1 + 1)) / (TF + k1 * (1 - b + b * DL / AVGDL)))
        # IDF = ln((N - n + 0.5) / (n + 0.5) + 1)
        
        sql = text("""
        WITH q_terms AS (
            SELECT id, doc_frequency 
            FROM terms 
            WHERE term = ANY(:tokens)
        ),
        doc_matches AS (
            SELECT 
                p.document_id,
                t.doc_frequency,
                p.term_frequency,
                d.doc_length
            FROM postings p
            JOIN q_terms t ON p.term_id = t.id
            JOIN documents d ON p.document_id = d.id
        )
        SELECT 
            document_id,
            SUM(
                LN((:total_docs - doc_frequency + 0.5) / 
                   (doc_frequency + 0.5) + 1) *
                ((term_frequency * (:k1 + 1)) / 
                 (term_frequency + :k1 * 
                  (1 - :b + :b * (doc_length / :avg_doc_length))))
            ) as score
        FROM doc_matches
        GROUP BY document_id
        ORDER BY score DESC
        LIMIT :limit
        """)
        
        result = await session.execute(sql, {
            "tokens": token_list,
            "total_docs": total_docs,
            "avg_doc_length": float(avg_doc_length),
            "k1": self.k1,
            "b": self.b,
            "limit": top_k
        })
        
        return [
            {"document_id": row.document_id, "score": float(row.score)}
            for row in result
        ]
\end{lstlisting}

This implementation calculates BM25 scores entirely within PostgreSQL using Common Table Expressions (CTEs). The query first identifies relevant terms from the query tokens (q\_terms), joins with the postings table to retrieve term frequencies and document lengths (doc\_matches), then computes the final BM25 score using PostgreSQL's mathematical functions. This approach achieves sub-second query latency even with large indices by leveraging PostgreSQL's query optimizer and the composite index on (term\_id, document\_id).

\clearpage

\section{Vector Search with HNSW}

\subsection{Embedding Generation Service}

The embedding service uses sentence-transformers with CLIP for multimodal embeddings, generating 512-dimensional dense vectors for both text and images. The service is implemented as a singleton to avoid repeated model loading.

\lstset{
  language=Python,
  basicstyle=\ttfamily\footnotesize,
  keywordstyle=\color{blue}\bfseries,
  commentstyle=\color{green!60!black}\itshape,
  stringstyle=\color{red},
  showstringspaces=false,
  breaklines=true,
  frame=single,
  numbers=left,
  numberstyle=\tiny\color{gray},
  caption={Embedding service for text and image encoding using CLIP},
  label=code:embedding_service
}
\begin{lstlisting}
class EmbeddingService:
    """Service for generating embeddings using CLIP model"""
    
    def __init__(self):
        self.device = "cuda" if torch.cuda.is_available() else "cpu"
        print(f"Loading embedding model on {self.device}...")
        self.model = SentenceTransformer(
            settings.EMBEDDING_MODEL,  # "clip-ViT-B-32"
            device=self.device
        )
    
    def encode_text(self, texts: Union[str, List[str]]) -> np.ndarray:
        """Encode text into 512-dim embeddings"""
        if isinstance(texts, str):
            texts = [texts]
        
        embeddings = self.model.encode(
            texts,
            convert_to_tensor=False,
            show_progress_bar=False,
            normalize_embeddings=True  # L2 normalization for cosine similarity
        )
        return embeddings
    
    def encode_image_base64(self, base64_string: str) -> np.ndarray:
        """Encode image from base64 string"""
        # Remove header if present (e.g. "data:image/jpeg;base64,")
        if "," in base64_string:
            base64_string = base64_string.split(",")[1]
        
        image_data = base64.b64decode(base64_string)
        image = Image.open(BytesIO(image_data))
        
        embedding = self.model.encode(
            image,
            convert_to_tensor=False,
            show_progress_bar=False,
            normalize_embeddings=True
        )
        return embedding
\end{lstlisting}

The encode\_image\_base64 method enables the system to process images uploaded by users through the web interface. Images are decoded from base64, converted to PIL Image objects, then encoded using CLIP's vision encoder. The normalize\_embeddings=True parameter ensures embeddings are L2-normalized, making cosine similarity equivalent to dot product and enabling faster similarity computation.

\subsection{HNSW Vector Search Implementation}

Vector search queries the pgvector HNSW index using PostgreSQL's custom operators. The \texttt{<=>} operator computes cosine distance efficiently using the HNSW graph structure.

\lstset{
  language=Python,
  basicstyle=\ttfamily\footnotesize,
  keywordstyle=\color{blue}\bfseries,
  commentstyle=\color{green!60!black}\itshape,
  stringstyle=\color{red},
  showstringspaces=false,
  breaklines=true,
  frame=single,
  numbers=left,
  numberstyle=\tiny\color{gray},
  caption={Vector similarity search using pgvector HNSW index},
  label=code:vector_search
}
\begin{lstlisting}
async def _vector_search(
    self,
    query_embedding: np.ndarray,
    session: AsyncSession,
    limit: int
) -> Dict[int, float]:
    """Perform vector similarity search using HNSW index"""
    try:
        # Convert numpy array to PostgreSQL vector format
        embedding_list = query_embedding.tolist()
        embedding_str = str(embedding_list)
        
        # Use pgvector's cosine similarity operator (<=>)
        # Returns documents ordered by ascending cosine distance
        query = text(f"""
            SELECT 
                de.document_id,
                1 - (de.embedding <=> '{embedding_str}'::vector) as similarity
            FROM document_embeddings de
            ORDER BY de.embedding <=> '{embedding_str}'::vector
            LIMIT :limit
        """)
        
        result = await session.execute(query, {"limit": limit})
        
        # Return as dict: document_id -> similarity score
        return {row.document_id: row.similarity for row in result}
        
    except Exception as e:
        print(f"Vector search error: {str(e)}")
        return {}
\end{lstlisting}

The \texttt{<=>} operator leverages the HNSW index (created with m=16, ef\_construction=64) to perform approximate nearest neighbor search in sublinear time. The HNSW graph structure enables efficient graph traversal, typically visiting only O(log N) nodes for a K-NN query rather than the O(N) required for exhaustive search.

\clearpage

\section{Hybrid Search and Reciprocal Rank Fusion}

\subsection{Search Service Integration}

The SearchService class orchestrates the hybrid retrieval pipeline, combining BM25 and vector search results using weighted score fusion.

\lstset{
  language=Python,
  basicstyle=\ttfamily\footnotesize,
  keywordstyle=\color{blue}\bfseries,
  commentstyle=\color{green!60!black}\itshape,
  stringstyle=\color{red},
  showstringspaces=false,
  breaklines=true,
  frame=single,
  numbers=left,
  numberstyle=\tiny\color{gray},
  caption={Hybrid search combining BM25 and vector retrieval},
  label=code:hybrid_search
}
\begin{lstlisting}
async def search(
    self,
    query: str,
    session: AsyncSession,
    top_k: int = None
) -> List[Dict[str, Any]]:
    """Hybrid search combining BM25 and vector similarity"""
    
    if top_k is None:
        top_k = settings.TOP_K_RESULTS  # Default: 10
    
    # Tokenize query for BM25
    tokenized_query = self.preprocess_query(query)
    
    # Generate query embedding for vector search
    query_embedding = self.embedding_service.encode_text(query)[0]
    
    # 1. Vector search using HNSW index (retrieve 2*top_k candidates)
    vector_results = await self._vector_search(
        query_embedding, session, top_k * 2
    )
    
    # 2. BM25 full-text search (retrieve 2*top_k candidates)
    bm25_results = await self._bm25_search(
        tokenized_query, session, top_k * 2
    )
    
    # 3. Combine and re-rank results
    combined_results = self._hybrid_rerank(
        vector_results, bm25_results, top_k
    )
    
    return combined_results

def _hybrid_rerank(
    self,
    vector_results: Dict[int, float],
    bm25_results: Dict[int, float],
    top_k: int
) -> List[Dict[str, Any]]:
    """Combine and re-rank results using weighted score fusion"""
    
    # Normalize scores to [0, 1] range
    def normalize_scores(scores: Dict[int, float]) -> Dict[int, float]:
        if not scores:
            return {}
        max_score = max(scores.values())
        min_score = min(scores.values())
        if max_score == min_score:
            return {k: 1.0 for k in scores}
        return {
            k: (v - min_score) / (max_score - min_score)
            for k, v in scores.items()
        }
    
    norm_vector = normalize_scores(vector_results)
    norm_bm25 = normalize_scores(bm25_results)
    
    # Combine scores with configurable weights
    # Default: VECTOR_WEIGHT=0.6, BM25_WEIGHT=0.4
    all_doc_ids = set(norm_vector.keys()) | set(norm_bm25.keys())
    combined_scores = {}
    
    for doc_id in all_doc_ids:
        vector_score = norm_vector.get(doc_id, 0.0)
        bm25_score = norm_bm25.get(doc_id, 0.0)
        
        # Weighted combination
        combined_score = (
            settings.VECTOR_WEIGHT * vector_score +
            settings.BM25_WEIGHT * bm25_score
        )
        combined_scores[doc_id] = combined_score
    
    # Sort by combined score descending
    sorted_results = sorted(
        combined_scores.items(),
        key=lambda x: x[1],
        reverse=True
    )[:top_k]
    
    return [
        {"document_id": doc_id, "score": score}
        for doc_id, score in sorted_results
    ]
\end{lstlisting}

The hybrid reranking algorithm first normalizes scores from both retrieval methods to a common [0,1] scale using min-max normalization. This ensures that BM25 scores (which can be arbitrarily large) and cosine similarities (bounded in [0,1]) are comparable. The weighted combination then applies configurable weights (default: 60\% vector, 40\% BM25) to produce the final ranking. This approach consistently outperforms either method alone, as vector search captures semantic similarity while BM25 provides precise lexical matching for technical terms and exact phrases.

\clearpage

\section{Distributed Web Crawling}

\subsection{Multi-threaded Crawler with DrissionPage}

The web crawler implements a multi-threaded architecture using DrissionPage for browser automation. The system supports both headless browser mode (for JavaScript-heavy sites) and fast session mode (for static content).

\lstset{
  language=Python,
  basicstyle=\ttfamily\footnotesize,
  keywordstyle=\color{blue}\bfseries,
  commentstyle=\color{green!60!black}\itshape,
  stringstyle=\color{red},
  showstringspaces=false,
  breaklines=true,
  frame=single,
  numbers=left,
  numberstyle=\tiny\color{gray},
  caption={Crawler worker implementation with multi-tab browser concurrency},
  label=code:crawler_worker
}
\begin{lstlisting}
class CrawlerWorker:
    """Crawler worker thread using DrissionPage"""
    
    def __init__(self, worker_id: int, stop_event: Event, browser_instance=None):
        self.worker_id = worker_id
        self.stop_event = stop_event
        self.url_manager = URLManager()
        self.content_extractor = ContentExtractor()
        self.page = None
        self.mode = DRISSION_MODE  # 's' for Session, 'd' for Driver
        self.browser_instance = browser_instance  # Shared browser for multi-tab
    
    def _create_page(self):
        """Create DrissionPage page object"""
        if self.mode == 's':
            # Session mode (fast, no JS execution)
            self.page = SessionPage(timeout=REQUEST_TIMEOUT)
            self.page.set.user_agent(USER_AGENT)
        else:
            # Driver mode (browser, JS execution)
            if self.browser_instance:
                # Multi-tab mode: create new tab in shared browser
                self.page = self.browser_instance.new_tab()
                self.page.set.timeouts(
                    base=REQUEST_TIMEOUT, 
                    page_load=REQUEST_TIMEOUT
                )
                logger.info(f"Worker {self.worker_id}: Created new tab")
            else:
                # Fallback: standalone browser instance
                co = ChromiumOptions()
                if HEADLESS:
                    co.headless()
                co.set_argument('--no-sandbox')
                co.set_argument('--disable-dev-shm-usage')
                co.set_user_agent(USER_AGENT)
                
                self.page = ChromiumPage(addr_or_opts=co)
                self.page.set.timeouts(
                    base=REQUEST_TIMEOUT, 
                    page_load=REQUEST_TIMEOUT
                )
    
    def fetch_page(self, url: str) -> Optional[str]:
        """Fetch page content with retry logic"""
        if not self.page:
            self._create_page()
        
        retry_count = 0
        while retry_count < MAX_RETRIES:
            try:
                # Navigate to URL
                self.page.get(url)
                
                # Wait for JS rendering if in browser mode
                if self.mode == 'd':
                    time.sleep(3)  # Wait for JS to complete
                
                # Extract HTML
                html = self.page.html
                
                if not html or len(html) < 100:
                    logger.warning(f"Empty HTML from {url}")
                    return None
                
                return html
            
            except (ElementNotFoundError, PageDisconnectedError) as e:
                logger.warning(f"Page error for {url}: {e}")
                retry_count += 1
                if retry_count < MAX_RETRIES:
                    time.sleep(2 ** retry_count)  # Exponential backoff
                    # Recreate page on error
                    self.page = None
            except Exception as e:
                logger.warning(f"Failed to fetch {url}: {e}")
                retry_count += 1
                if retry_count < MAX_RETRIES:
                    time.sleep(2 ** retry_count)
        
        return None
\end{lstlisting}

The crawler worker implements exponential backoff for failed requests, doubling the wait time on each retry (2, 4, 8 seconds). The multi-tab architecture allows multiple workers to share a single browser instance while crawling different URLs concurrently, significantly reducing memory footprint compared to launching separate browser processes per worker.

\subsection{URL Management with Redis and Bloom Filter}

The URL manager uses Redis for distributed URL queue management and a Bloom filter for efficient duplicate URL detection across distributed workers.

\lstset{
  language=Python,
  basicstyle=\ttfamily\footnotesize,
  keywordstyle=\color{blue}\bfseries,
  commentstyle=\color{green!60!black}\itshape,
  stringstyle=\color{red},
  showstringspaces=false,
  breaklines=true,
  frame=single,
  numbers=left,
  numberstyle=\tiny\color{gray},
  caption={URL management with Redis queue and Bloom filter deduplication},
  label=code:url_manager
}
\begin{lstlisting}
class URLManager:
    """Distributed URL queue manager using Redis"""
    
    def __init__(self):
        self.redis_client = redis.Redis(
            host=REDIS_HOST,
            port=REDIS_PORT,
            decode_responses=True
        )
        self.bloom = BloomFilter(max_elements=10000000, error_rate=0.001)
        self._load_bloom_filter()
    
    def add_urls(self, urls: List[str], depth: int = 0) -> int:
        """Add URLs to queue if not already visited"""
        added = 0
        for url in urls:
            if not self.is_visited(url):
                # Push to Redis queue with depth metadata
                self.redis_client.rpush(
                    "crawler:url_queue",
                    json.dumps({"url": url, "depth": depth})
                )
                # Add to Bloom filter for fast duplicate detection
                self.bloom.add(url)
                added += 1
        return added
    
    def get_next_url(self) -> Optional[dict]:
        """Pop next URL from queue (FIFO)"""
        item = self.redis_client.lpop("crawler:url_queue")
        if item:
            return json.loads(item)
        return None
    
    def is_visited(self, url: str) -> bool:
        """Check if URL already crawled using Bloom filter"""
        # Bloom filter provides probabilistic membership test
        # False positives possible, but no false negatives
        return url in self.bloom
    
    def mark_visited(self, url: str):
        """Mark URL as visited"""
        self.bloom.add(url)
        self.redis_client.sadd("crawler:visited_urls", url)
    
    def check_and_rest_if_needed(self):
        """Rate limiting: enforce politeness delay"""
        current_time = time.time()
        last_request = float(self.redis_client.get("crawler:last_request") or 0)
        
        # Ensure minimum delay between requests (default: 1 second)
        if current_time - last_request < REQUEST_DELAY:
            time.sleep(REQUEST_DELAY - (current_time - last_request))
        
        self.redis_client.set("crawler:last_request", current_time)
\end{lstlisting}

The Bloom filter uses 10 million maximum elements with 0.1\% error rate, requiring approximately 18MB of memory. This enables the crawler to efficiently check whether a URL has been visited without querying Redis for every URL, reducing network latency. The false positive rate of 0.1\% means only 0.1\% of new URLs might be incorrectly skipped, which is acceptable for web crawling applications.

\clearpage

\section{Image Extraction and Multimodal Indexing}

\subsection{Image Extraction from Web Pages}

The image extractor identifies and downloads images from crawled pages, filtering by size and format to retain only meaningful visual content.

\lstset{
  language=Python,
  basicstyle=\ttfamily\footnotesize,
  keywordstyle=\color{blue}\bfseries,
  commentstyle=\color{green!60!black}\itshape,
  stringstyle=\color{red},
  showstringspaces=false,
  breaklines=true,
  frame=single,
  numbers=left,
  numberstyle=\tiny\color{gray},
  caption={Image extraction and processing for multimodal search},
  label=code:image_extraction
}
\begin{lstlisting}
class ImageExtractor:
    """Extract and process images from HTML"""
    
    MIN_WIDTH = 200
    MIN_HEIGHT = 200
    MAX_IMAGE_SIZE = 5 * 1024 * 1024  # 5MB
    
    def extract_images(self, html: str, base_url: str) -> List[Dict]:
        """Extract images with metadata from HTML"""
        soup = BeautifulSoup(html, 'html.parser')
        images = []
        
        for img_tag in soup.find_all('img'):
            try:
                # Get image source
                img_url = img_tag.get('src') or img_tag.get('data-src')
                if not img_url:
                    continue
                
                # Resolve relative URLs
                absolute_url = urljoin(base_url, img_url)
                
                # Download image
                response = requests.get(
                    absolute_url, 
                    timeout=10,
                    headers={'User-Agent': USER_AGENT}
                )
                
                if response.status_code != 200:
                    continue
                
                # Check file size
                if len(response.content) > self.MAX_IMAGE_SIZE:
                    logger.warning(f"Image too large: {absolute_url}")
                    continue
                
                # Open and validate image
                img = Image.open(BytesIO(response.content))
                width, height = img.size
                
                # Filter small images (likely icons/logos)
                if width < self.MIN_WIDTH or height < self.MIN_HEIGHT:
                    continue
                
                # Convert to base64 for storage
                buffered = BytesIO()
                img.save(buffered, format=img.format or 'PNG')
                img_base64 = base64.b64encode(
                    buffered.getvalue()
                ).decode('utf-8')
                
                # Extract metadata
                alt_text = img_tag.get('alt', '')
                title = img_tag.get('title', '')
                
                images.append({
                    'url': absolute_url,
                    'base64': img_base64,
                    'width': width,
                    'height': height,
                    'alt_text': alt_text,
                    'title': title,
                    'format': img.format
                })
                
            except Exception as e:
                logger.warning(f"Failed to extract image: {e}")
                continue
        
        logger.info(f"Extracted {len(images)} images from {base_url}")
        return images
\end{lstlisting}

The extractor applies multiple quality filters: (1) minimum dimensions of 200x200 pixels to exclude icons and thumbnails, (2) maximum file size of 5MB to avoid extremely large images, and (3) format validation to ensure the image is readable. Extracted images are converted to base64 encoding for storage in the database, eliminating the need for separate file system management.

\clearpage

\section{HTTP Microservice for Embedding Generation}

\subsection{FastAPI Service for Image Encoding}

A dedicated Python microservice handles image embedding generation, exposing a RESTful API that accepts base64-encoded images and returns CLIP embeddings.

\lstset{
  language=Python,
  basicstyle=\ttfamily\footnotesize,
  keywordstyle=\color{blue}\bfseries,
  commentstyle=\color{green!60!black}\itshape,
  stringstyle=\color{red},
  showstringspaces=false,
  breaklines=true,
  frame=single,
  numbers=left,
  numberstyle=\tiny\color{gray},
  caption={FastAPI microservice for image embedding generation},
  label=code:image_service
}
\begin{lstlisting}
from fastapi import FastAPI, HTTPException
from pydantic import BaseModel
from embedding_service import get_embedding_service
import uvicorn

app = FastAPI(title="Verdant Search Embedding Service")
embedding_service = None

class ImageRequest(BaseModel):
    image_base64: str
    
class ImageResponse(BaseModel):
    embedding: List[float]
    dimension: int

@app.on_event("startup")
async def startup_event():
    """Initialize embedding service on startup"""
    global embedding_service
    embedding_service = get_embedding_service()
    print("Embedding service initialized")

@app.post("/encode/image", response_model=ImageResponse)
async def encode_image(request: ImageRequest):
    """
    Encode image to CLIP embedding vector
    
    Args:
        image_base64: Base64-encoded image string
        
    Returns:
        512-dimensional CLIP embedding vector
    """
    try:
        # Generate embedding
        embedding = embedding_service.encode_image_base64(
            request.image_base64
        )
        
        return ImageResponse(
            embedding=embedding.tolist(),
            dimension=len(embedding)
        )
        
    except Exception as e:
        raise HTTPException(
            status_code=500,
            detail=f"Failed to encode image: {str(e)}"
        )

@app.post("/encode/text", response_model=ImageResponse)
async def encode_text(text: str):
    """Encode text to embedding vector"""
    try:
        embedding = embedding_service.encode_text(text)[0]
        
        return ImageResponse(
            embedding=embedding.tolist(),
            dimension=len(embedding)
        )
        
    except Exception as e:
        raise HTTPException(
            status_code=500,
            detail=f"Failed to encode text: {str(e)}"
        )

@app.get("/health")
async def health_check():
    """Health check endpoint"""
    return {
        "status": "healthy",
        "model": settings.EMBEDDING_MODEL,
        "device": embedding_service.device if embedding_service else "N/A"
    }

if __name__ == "__main__":
    uvicorn.run(app, host="0.0.0.0", port=8001)
\end{lstlisting}

The microservice architecture isolates the GPU-dependent embedding generation from other system components. This enables independent scaling: multiple embedding service instances can run on GPU-equipped nodes while the database and search services run on CPU-only nodes. The health check endpoint reports the loaded model and device (CPU/CUDA), enabling monitoring systems to verify GPU utilization.

\clearpage

\section{Performance Optimizations}

\subsection{Redis Caching Layer}

To reduce database load and improve query latency, the system implements a Redis caching layer that stores search results and embeddings for frequently accessed queries.

\lstset{
  language=Python,
  basicstyle=\ttfamily\footnotesize,
  keywordstyle=\color{blue}\bfseries,
  commentstyle=\color{green!60!black}\itshape,
  stringstyle=\color{red},
  showstringspaces=false,
  breaklines=true,
  frame=single,
  numbers=left,
  numberstyle=\tiny\color{gray},
  caption={Redis caching for search results and query embeddings},
  label=code:redis_cache
}
\begin{lstlisting}
class CachedSearchService:
    """Search service with Redis caching"""
    
    def __init__(self):
        self.search_service = SearchService()
        self.redis = redis.Redis(
            host=settings.REDIS_HOST,
            decode_responses=False  # Store binary data (embeddings)
        )
        self.CACHE_TTL = 300  # 5 minutes
    
    async def search_with_cache(
        self,
        query: str,
        session: AsyncSession,
        top_k: int = 10
    ) -> List[Dict[str, Any]]:
        """Search with result caching"""
        
        # Generate cache key from query and parameters
        cache_key = f"search:{hash(query)}:{top_k}"
        
        # Try to get from cache
        cached = self.redis.get(cache_key)
        if cached:
            logger.info(f"Cache hit for query: {query}")
            return json.loads(cached)
        
        # Cache miss: execute search
        results = await self.search_service.search(
            query, session, top_k
        )
        
        # Store in cache with TTL
        self.redis.setex(
            cache_key,
            self.CACHE_TTL,
            json.dumps(results)
        )
        
        return results
    
    def cache_embedding(self, text: str, embedding: np.ndarray):
        """Cache text embedding to avoid recomputation"""
        cache_key = f"embedding:{hash(text)}"
        
        # Store numpy array as bytes
        embedding_bytes = embedding.tobytes()
        self.redis.setex(
            cache_key,
            self.CACHE_TTL * 2,  # Longer TTL for embeddings
            embedding_bytes
        )
    
    def get_cached_embedding(self, text: str) -> Optional[np.ndarray]:
        """Retrieve cached embedding"""
        cache_key = f"embedding:{hash(text)}"
        
        cached = self.redis.get(cache_key)
        if cached:
            # Reconstruct numpy array from bytes
            return np.frombuffer(cached, dtype=np.float32)
        
        return None
\end{lstlisting}

The caching strategy differentiates between complete search results (5-minute TTL) and intermediate embeddings (10-minute TTL). Embeddings have a longer TTL because they are more expensive to compute (requiring GPU inference) and are reusable across different search queries. Cache keys use Python's hash() function for fast key generation, accepting potential hash collisions given the short TTL and probabilistic nature of caching.

\subsection{Batch Processing for Index Updates}

Index updates are batched to reduce database transaction overhead and improve throughput during bulk document ingestion.

\lstset{
  language=Python,
  basicstyle=\ttfamily\footnotesize,
  keywordstyle=\color{blue}\bfseries,
  commentstyle=\color{green!60!black}\itshape,
  stringstyle=\color{red},
  showstringspaces=false,
  breaklines=true,
  frame=single,
  numbers=left,
  numberstyle=\tiny\color{gray},
  caption={Batch indexing for improved throughput},
  label=code:batch_index
}
\begin{lstlisting}
class BatchIndexService:
    """Service for batched document indexing"""
    
    BATCH_SIZE = 100
    
    def __init__(self):
        self.batch_buffer = []
        self.lock = asyncio.Lock()
    
    async def add_to_batch(
        self,
        doc: Dict[str, Any],
        session: AsyncSession
    ):
        """Add document to batch, flush if batch full"""
        async with self.lock:
            self.batch_buffer.append(doc)
            
            if len(self.batch_buffer) >= self.BATCH_SIZE:
                await self._flush_batch(session)
    
    async def _flush_batch(self, session: AsyncSession):
        """Flush accumulated batch to database"""
        if not self.batch_buffer:
            return
        
        logger.info(f"Flushing batch of {len(self.batch_buffer)} documents")
        
        # 1. Bulk insert documents
        doc_ids = await self._bulk_insert_documents(
            self.batch_buffer, session
        )
        
        # 2. Batch generate embeddings (GPU efficient)
        texts = [doc['content'] for doc in self.batch_buffer]
        embeddings = self.embedding_service.encode_text(texts)
        
        # 3. Bulk insert embeddings
        await self._bulk_insert_embeddings(
            doc_ids, embeddings, session
        )
        
        # 4. Update inverted index in batch
        await self._batch_update_inverted_index(
            doc_ids, self.batch_buffer, session
        )
        
        # Clear buffer
        self.batch_buffer = []
        
        logger.info(f"Batch flush completed")
    
    async def _bulk_insert_documents(
        self,
        docs: List[Dict],
        session: AsyncSession
    ) -> List[int]:
        """Bulk insert using COPY or INSERT INTO VALUES"""
        # Use PostgreSQL COPY for maximum throughput
        stmt = text("""
            INSERT INTO documents (title, content, url, source_type)
            VALUES (:title, :content, :url, :source_type)
            RETURNING id
        """)
        
        results = []
        for doc in docs:
            result = await session.execute(stmt, doc)
            results.append(result.scalar())
        
        await session.commit()
        return results
\end{lstlisting}

Batching reduces the number of database round-trips by 100x (for BATCH\_SIZE=100), dramatically improving indexing throughput. The GPU embedding generation also benefits from batching: encoding 100 texts in a single batch is 10-20x faster than encoding them individually due to reduced model invocation overhead and better GPU utilization through parallel matrix operations.

\clearpage
